\documentclass{article}
\usepackage{graphicx} % Required for inserting images

\title{Chapter 23 Problem Set}
\author{Siddharth K}
\date{April 2024}

\begin{document}

\maketitle

\textbf{1a.} Let's say the first countable sets is $A$ and $B$ is the second countable set excluding the intersection of the two sets. Also, let's say set $C$ contains the first element of $A$, the first element of $B$, the second element of $A$, the second element of $B$, and so on in that order.
\newline

Since $A$ is countable, there is a function with an inverse that maps $A$ to $N$. If $a \in A$ maps to $n \in N$, then let $a \in C$ map to $2n - 1 \in N$. Likewise, if $b\in B$ maps to $n \in N$, let $b \in C$ map to $2n \in N$. This function maps every value in $A \cup B$ to $N$ and it has an inverse.
\newline

Now, let's consider the sets $A_1, A_2,\ldots, A_n$, and let's say that $B = \bigcup_{i = 1}^n A_i$. Let's assume $B$ is countable. Then, by the logic from earlier, if $A$ is countable, then $A \cup B$ is countable. By induction, the union of a finite number of countable sets is also countable.
\newline

\textbf{1b.} Let's say the countable sets are $A_1, A_2, \ldots$. We can arrange the elements of these sets diagonally. For example, in the top row, we can write the first element of $A_1$. In the second row, we can write the first element of $A_2$ and the second element of $A_1$. In the third row, we can write the first element of $A_3$, second element of $A_2$, and the third element of $A_3$. In general, the $k$'th number in the $n$'th row is the $k$'th element of $A_{n - k + 1}$. We can assign a natural number to each of these numbers in the grid, so we have established a bijection to $N$.
\newline

\textbf{2a.} Assume that none of the $A_i$ are infinite. Then they all have a finite cardinality; we can call this $n_i$. That means the cardinality of $\bigcup_{i = 1}^n A_i$ is at most $\sum_{i = 1}^n n_i$, which is a finite number. Therefore, at least one of $A_i$ must be infinite.
\newline

\textbf{2b.} Assume all $A_i$ is countable. By problem 1a, their union is also countable, which is a contradiction. So at least one of them must be uncountable.
\newline

\textbf{3a.} Assume $f(s) = t \in T$. Then $g(f(s)) = g(t) = s$. Since $g$ is a function from $T$ to $S$, each value in $T$ only corresponds to one in $S$. Equivalently, if $s, s' \in S$ and $s \neq s'$ then $f(s) \neq f(s')$, which is the definition of injectivity.
\newline

Going the other way, we start with the fact that $f(s) \neq f(s')$ for all $s, s' \in S$ where $s \neq s'$. Now, let's consider a function $g: T \rightarrow S$. We know that no two values in $S$ correspond to the same value in $T$, so $g(t) = s$ for some unique $s \in S$. Since $t = f(s)$, we have $g(f(s)) = s$.
\newline

\textbf{3b.} Let's say $g(t) = s \in S$. We know that $s$ exists because $g$ maps all elements in $T$ to some in $S$. So $f(g(t)) = f(s) = t$. This tells us that for all $t \in T$, there exists $s \in S$ such that $f(s) = t$, which is the definition of surjectivity.
\newline

Now, let's prove the converse. We assume that for all $t \in T$, there exists $s \in S$ such that $f(s) = t$. So every element in $T$ has a corresponding value in $S$, which means $f$ has an inverse, which we can call $g$. Therefore, $f(g(t)) = t$.
\newline

\textbf{4.} Since $Q$ is countable, we can assign a unique natural number to every rational number and all of them will be covered. Therefore, it would suffice to prove that $N \times N$ is countable. To do this, let's plot all ordered pairs $(a, b)$ where $a, b \in N$ on a plane. Then, in a similar manner to the proof that $Q$ is countable, we can form arrows starting at $(1, 1)$, then moving to $(2, 1)$, then $(1, 2)$, then $(1, 3)$, then $(2, 2)$, and so on in that pattern. Each point can be represented by the number of arrows drawn to reach it. This is a bijection from $N$ to $N \times N$.
\newline

\textbf{5.} First, we can map a number in $(0, 1)$ to another in $(1, \infty)$ by taking its reciprocal. Then, we can map a number in $(1, \infty)$ to $(0, \infty)$ by subtracting $1$. Now, we can map a number $n \in (0, \infty)$ to $\log(n) \in (-\infty, \infty)$.
\newline

\textbf{6.} When trying to create a new rational number by making each digit different from the same digit of at least one of the rational numbers listed, there is no guarantee that the resulting number is also rational.
\newline

\textbf{7.} By problem 1a, the union of two countable sets is also countable. That means, if $A\backslash B$ is countable, then $A \cup A\backslash B = A$ is uncountable, which is a contradiction. Therefore, $A\backslash B$ is uncountable.
\newline

Let the elements of $B$ be $b_0, b_1, b_2, \ldots$ and a countable subset of $A$ have elements $a_0, a_1, a_2, \ldots$. Then we can map each $a_i$ to $b_i$ for all $i \ge 0$. Then, for the elements in $A$ that are not in $B$, we can map them to themselves. This is a bijection, so $A$ and $A \backslash B$ have the same cardinality.
\newline

\textbf{8.} Consider a function $f$ that maps a binary string $s$ to a real number in the interval $[0, 1)$ by writing the string after a decimal point, and then converting it to a decimal number. For example, the string $101$ would correspond to the number $0.101_2 = 0.625$. The function $f$ is injective because no two binary numbers have the same decimal expansion, and the string $111...$ corresponds to $0.\overline{9}$, which is the maximum number in $[0, 1)$. Also, $f$ is surjective because no decimal number in $[0, 1)$ has two different binary representations. So $f$ is a bijection between infinite binary strings and real numbers in $[0, 1)$.
\newline

\textbf{9.} I tried to find a function that maps specific binary strings to ternary strings, but that couldn't work since the number of finite binary strings is a power of $2$ and the number of finite ternary strings is a power of $3$, and those two can never be equal. I couldn't figure out how else to solve the problem.

\end{document}